\uuid{G01F}
\exo7id{3546}
\titre{exo7 3546}
\auteur{quercia}
\organisation{exo7}
\datecreate{2010-03-10}
\isIndication{false}
\isCorrection{true}
\chapitre{Réduction d'endomorphisme, polynôme annulateur}
\sousChapitre{Polynôme annulateur}
\cpge{mpsi-mp2i.calcul-matriciel-systemes.e, pcsi.determinants.c, ptsi.determinants.a, mpsi-mp2i.matrices-applications-lineaires.b, pcsi.matrices-applications-lineaires.b, ptsi.matrices-applications-lineaires.b, bcpst1.matrices, mp-mpi.reduction-endomorphismes.d, pc.reduction.c, psi.reduction-endomorphismes.c, pt.reduction-endomorphismes.c, mp-mpi.reduction-endomorphismes.g, pc.complements-espaces-vectoriels.d, psi.complements-espaces-vectoriels.d, mp-mpi.equations-differentielles-lineaires.c, bcpst2.valeurs-propres-vecteurs-propres.b}
\module{Algèbre}
\niveau{L2}
\difficulte{}

\contenu{
\question{
Soit $A \in \mathcal{M}_n(\C)$ telle que $A = A^{-1}$. $A$ est-elle diagonalisable ?
Calculer $e^A$. ($e^A = \sum_{k=0}^\infty \frac{A^k}{k!}$)
}
\reponse{
$A$ est diagonalisable car $A^2=I$. $e^A = (\ch1)I + (\sh1)A$.
}
}
